% Final width: 110 mm. Inherits the volume palette and figure font.
\begin{tikzpicture}[x=1mm,y=1mm,
  every node/.style={font=\figurefont\footnotesize,text=FigureInk},
  box/.style={draw=FigureStroke,rounded corners=1mm,align=center,
    minimum height=16mm,inner sep=2mm},
  flow/.style={->,draw=FigureStroke,line width=.7pt}]
  \node[box,text width=25mm,draw=FigureBlue] (omega) at (15,30)
    {基础概率空间\\\((\Omega,\mathbb P)\)};
  \node[box,text width=26mm,draw=LancetGreen] (value) at (54,30)
    {整个读数规则\\\(X:\Omega\to\mathbb R\)};
  \node[box,text width=27mm,draw=FigureRed] (law) at (94,30)
    {取值的分布\\\(P_X(\{9\})=1/2\)};
  \draw[flow] (omega) -- (value);
  \draw[flow] (value) -- node[above=1mm] {推送} (law);
  \node[align=center,text=FigureMuted] at (54,16)
    {单次观测：\(\omega=b\mapsto X(b)=9\)；推送作用于整个概率规律};
  \node[box,text width=99mm,minimum height=13mm] at (54,1)
    {规律未知：样本\((x_1,\ldots,x_n)\) \(\longrightarrow\)
     估计目标\(\mathbb E X\)、分位数或分布\\
     同一分布可以由不同基础空间上的映射产生};
\end{tikzpicture}
